\scanpage{83}
% loc 1–4, struck block
\deleted{deux \deleted{éléments} parmi les éléments $1,u,v,uv$ est dans ce sous-module. Il suffit de \deleted{le} voir pour deux tels éléments distincts de 1. \deleted{\unsure{Il|On}} faut regarder}
% loc 5–7, struck multiplication table
\[\begin{aligned}
\deleted{u.u}&\deleted{=\det(u)1+(\operatorname{Tr}u)u+\unsure{\det}}\\
\deleted{u.v}&\deleted{=0.1+0.u+0.v+1uv}\\
\deleted{u.uv}&\deleted{=}
\end{aligned}\]
% loc 8, crossed list
\[ \deleted{\underline{u^2},\underline{v^2},uvuv,\underline{uv},u^2v,\underline{vu},vuv,uvu,uv^2} \]
% loc 9–12, struck calculations
\deleted{Or \unsure{$u^2$}}
\[\begin{aligned}
\deleted{u^2}&\deleted{=\operatorname{Tr}(u).u-\det(u).1}&&\deleted{\text{((2) pour }u}\\
\deleted{v^2}&\deleted{=\operatorname{Tr}(v).v.-\det(v)1}&&\deleted{\text{(2) pour }v}\\
\deleted{uv}&\deleted{=1.uv}&&\deleted{\text{tautol}}\\
\deleted{vu}&\deleted{=-1.uv+\operatorname{Tr}(v)u+\operatorname{Tr}u\,v+(\operatorname{Tr}(uv)-\operatorname{Tr}u\operatorname{Tr}v)1}
\end{aligned}\]
\deleted{(relation $(6')$)}
% loc 13–18, new proof
Ce sous-module $V$ est stable par multiplication \inserted{\uline{à gauche}} par \unsure{$1$}, $u,v$, \deleted{$uv$ \unsure{.} \illegible{3}} \inserted{(donc aussi par} $uv$, sans compter 1). Donc il suffit de voir que $ux,vx\in V$, si $x\in V$. Il suffit de prendre pour $x$ un des \unsure{gén|gèn} $1,u,v,uv$. Le cas 1 \unsure{étant|est} trivial, il reste
% loc 19
\[ \underline{u^2},\underline{uv},\underline{u^2}v,\;vu,\underline{v^2},vuv=\deleted{\unsure{v(uv)}}(vu)v \]
% loc 20–23
Pour $u^2,v^2$ on utilise H.C. (3), pour $uv$ c'est tautologique, pour $u^2v$ cela résulte de l'expression de $u^2$ H.C., pour $vu$ on utilise $(6')$, qu'on écrit
% loc 24, overwritten formula
\[ \deleted{vu=-\unsure{uv}+\operatorname{Tr}(v).u+{}} \]
% loc 25
\[ vu=(\operatorname{Tr}(uv)-\operatorname{Tr}u,\operatorname{Tr}v)1+\operatorname{Tr}(v)u+\operatorname{Tr}(u)v-uv. \tag*{$(6'')$} \]
% loc 26–29
Enfin pour \deleted{$(vu)v$} $vuv=(vu)v$, on multiplie la relation $(6'')$ à droite par $v$, en utilisant la relation H.C. pour $v$ (pour se débarrasser des termes en $v^2$).
% archive sheet number: 82
